\chapter{User Stories \\
\small{\textit{-- Ryan Davis, Cory Vitanza, and Roy Tas}}
\index{glossary} 
\index{Chapter!User Stories}
\label{Chapter::UserStories}}

\section{Overview}
This chapter presents key user stories that describe the real-world needs and motivations of potential users of the tracking system. These stories help guide the system’s design and ensure that functional and non-functional requirements align with user expectations.

\section{User Stories}

\subsection{Parent of a Child with Autism}
\begin{itemize}
    \item \StoryLabel{usParentAutism}{US1}: \textbf{As a parent of a child with autism,} I want a device that automatically alerts me when my child leaves a safe zone so that I can respond quickly without wasting time searching.
\end{itemize}

\subsection{Caregiver in an Assisted Living Facility}
\begin{itemize}
    \item \StoryLabel{usCaregiver}{US2}: \textbf{As a caregiver responsible for multiple residents,} I want to track multiple devices simultaneously and receive alerts when a resident leaves a designated safe area so that I can act promptly to ensure their safety.
\end{itemize}

\subsection{Family Member Monitoring Battery Life}
\begin{itemize}
    \item \StoryLabel{usBattery}{US3}: \textbf{As a family member monitoring a loved one,} I want to receive notifications when the tracker’s battery is running low so that I can prevent the device from powering off unexpectedly during an emergency.
\end{itemize}

\subsection{User Triggering an Emergency Alert}
\begin{itemize}
    \item \StoryLabel{usEmergency}{US4}: \textbf{As a user wearing the device,} I want to press a button in an emergency to send an alert to my caregiver or family members so that I can get help quickly.
\end{itemize}

\subsection{System Administrator Managing Devices}
\begin{itemize}
    \item \StoryLabel{usSysAdmin}{US5}: \textbf{As a system administrator,} I want to manage LoRaWAN gateways and device registrations efficiently so that the system remains scalable and reliable without subscription costs.
\end{itemize}

\subsection{AI Monitoring and Pattern Detection}
\begin{itemize}
    \item \StoryLabel{usAIPattern}{US6}: \textbf{As a system backend operator,} I want the system to automatically detect unusual movement patterns so that potential emergencies can be identified proactively.
\end{itemize}

% ---------- Mapping Table ----------
\section{User Story to Use Case Mapping}

\begin{longtable}{|p{3cm}|p{7cm}|p{4.5cm}|}
\caption{User Story to Use Case Mapping \label{Table::UserStoryMapping}}\\
\hline
\textbf{User Story ID} & \textbf{Description} & \textbf{Mapped Use Case(s)} \\
\hline
\StoryReference{usParentAutism} & Parent receives alerts when a child leaves a safe zone. & \UseCaseReference{ucGeofenceAlert} \\
\hline
\StoryReference{usCaregiver} & Caregiver monitors multiple residents and receives alerts. & \UseCaseReference{ucMultiResidentMonitoring} \\
\hline
\StoryReference{usBattery} & Family member receives low-battery notifications. & \UseCaseReference{ucBatteryStatusAlert} \\
\hline
\StoryReference{usEmergency} & User triggers an emergency alert to notify caregiver. & \UseCaseReference{ucEmergencyAlert} \\
\hline
\StoryReference{usSysAdmin} & System administrator manages deployments and gateways. & \UseCaseReference{ucCostEfficientDeployment} \\
\hline
\StoryReference{usAIPattern} & Backend detects unusual behavior through ML models. & \UseCaseReference{ucPatternDetection} \\
\hline
\end{longtable}

\section{Traceability Discussion}
All user stories directly support at least one system use case to ensure end-to-end traceability. 
Each story reflects the needs of a real stakeholder—caregivers, users, family members, or system administrators—and guides the functional and non-functional requirements of the FindIt asset tracking system. 
By maintaining explicit links (\StoryReference{usParentAutism}–\StoryReference{usAIPattern}) between stories and use cases, this document ensures future system iterations remain aligned with user needs.
